\section{Conclusion}
\label{sec:conclusion}

This paper presented PhyST-Net, a physics-aware differentiable spatio-temporal Transformer for constrained network forecasting with historical observations, future exogenous drivers, graph metadata, and operational bounds. Rather than treating these inputs as a single fused tensor, PhyST-Net defines a structured forecasting contract: AGSP provides continuous temporal tokenization, R-STMF separates relational spatio-temporal mechanisms, K-AMC injects forecast-horizon drivers through non-linear adaptive modulation, and PI-DCH promotes feasible decoding through a differentiable constraint head.

The experimental design focuses the main leaderboard on wind-power forecasting, where physical capacity limits, geographic propagation, and meteorological forcing are all present. Across SDWPF, HoT, and Kelmarsh, PhyST-Net achieves the best MAE in all representative prediction-step settings and the best RMSE in most of them, with especially clear gains on Kelmarsh and SDWPF. Auxiliary ablations on ASHRAE building energy, HK Rooftop PV, and SDWPF further show that the proposed modules contribute differently across domains: relational manifold modeling, adaptive temporal patching, and future-driver conditioning become most important under different physical regimes.

Several directions remain open. The current graph interface is primarily static, while deployed infrastructure may experience topology changes, sensor outages, and maintenance-induced distribution shifts. The present formulation is also deterministic; extending PI-DCH to calibrated probabilistic forecasts under hard operational envelopes would better support risk-aware scheduling. Finally, larger-scale deployment should evaluate how the forecasting contract behaves under richer NWP products, additional energy assets, and online adaptation scenarios.
